\subsection*{Radius-matched first-edit timing}
Before the first incumbent edit, allow one opening from the first
two offers. The radius is a quarter of the median positive radius
at the first offer and is fixed thereafter. The ranker is shared
by the learned stop and earliest arms. The timing model predicts
the next-offer advantage from the current offer only. Its donor
labels use ranks fitted on other donors, and waiting can expire.

\begin{center}
\small
\setlength{\tabcolsep}{4pt}
\begin{tabular}{lrrrr}
\hline
Policy & Opens & Late & Risk & Joint \\
\hline
Incumbent & \OpeningIncumbent \\
Frozen rank, first & \OpeningFrozen \\
Learned rank, first & \OpeningFirst \\
Always wait & \OpeningWait \\
Learned stop & \OpeningLearned \\
Matched first & \OpeningMatched \\
\hline
\end{tabular}
\end{center}

On all 102 views, the matched first-offer control opens only where
the timing model opens; this is an offline utility control. The
paired contrast is \OpeningEffect{} and \OpeningExpired{} waits
expire in the actual policy. The original risk rule is unchanged.
The outcome-aware first and two-offer menus both cover
\OpeningBothJoint{} joint views; the learned top-one menu at
either offer covers only \OpeningTop. This is a candidate-ranking
failure, not evidence that a longer timing model would help.
